\documentclass[11pt]{article}
\usepackage[a4paper,margin=1in]{geometry}
\usepackage{amsmath,amsthm,amssymb,mathtools}
\usepackage{microtype}
\usepackage{hyperref}
\hypersetup{colorlinks=true,linkcolor=blue,citecolor=blue,urlcolor=blue}

\title{\textbf{Chebyshev--Lucas Error Correction in Decimated Transfer--Matrix Chains:\\
A Simple, First--Principles Framework with Robust Decoding}}
\author{VR \& AI}
\date{\today}

% theorem envs
\newtheorem{theorem}{Theorem}[section]
\newtheorem{proposition}[theorem]{Proposition}
\newtheorem{lemma}[theorem]{Lemma}
\theoremstyle{definition}
\newtheorem{definition}[theorem]{Definition}
\theoremstyle{remark}
\newtheorem{remark}[theorem]{Remark}
\newtheorem{example}[theorem]{Example}

% macros
\newcommand{\Z}{\mathbb{Z}}
\newcommand{\N}{\mathbb{N}}
\newcommand{\C}{\mathbb{C}}
\newcommand{\tr}{\mathrm{tr}}

\begin{document}
\maketitle

\begin{abstract}
We show that any one--dimensional medium governed by a constant $2\times2$ unit--cell transfer matrix $M$ admits a \emph{universal tri--site recurrence} when stroboscopically sampled every $l$ sites:
\[
x_{k+2l}\;-\;V_l\,x_{k+l}\;+\;Q^l\,x_k \;=\; 0,\qquad V_l=\tr(M^l),\ \ Q=\det M,\ \ Q^l=\det(M^l)=(\det M)^l.
\]
This identity follows directly from Cayley--Hamilton applied to the $l$--step monodromy $M^l$.
We recast it as a \emph{convolutional constraint} for \emph{event--coded} measurements and obtain a
\emph{Chebyshev--Lucas checksum} enabling noise--robust decoding via a finite--state trellis (Viterbi/dynamic programming).
We give rules for choosing the alphabet and window, prove existence of a unique maximum--likelihood (ML) decoded sequence on a finite window, derive logarithmic bounds on the occupied support radius under repeated injection, and extend the construction to weakly coupled multi--track systems by diagonalizing transverse modes.
We also clarify when ``symmetric regrouping'' into Dickson polynomials applies, provide a counterexample showing the checksum alone is not injective (hence normalization/ML decoding matters), and include concrete validation metrics.
\end{abstract}

\tableofcontents

\section{From unit cell to decimated tri--site identity}
\subsection{First principles: second order $\Rightarrow$ $2\times2$ transfer}
Many one--dimensional media reduce to a scalar second--order recurrence after eliminating a conjugate variable (momentum/flux):
\begin{equation}\label{eq:secondorder}
x_{k+1} \;=\; P\,x_k \;-\; Q\,x_{k-1},\qquad P,Q\in\C,\ Q\neq 0.
\end{equation}
Collecting into a first--order map,
\begin{equation}\label{eq:transfer}
\binom{x_{k+1}}{x_k}\;=\; M\,\binom{x_k}{x_{k-1}},\qquad
M=\begin{pmatrix} P & -Q\\[2pt] 1 & 0\end{pmatrix},\quad \tr M=P,\ \det M=Q.
\end{equation}

\subsection{Cayley--Hamilton on the $l$--step monodromy}
\begin{proposition}[Decimated recurrence]\label{prop:decimated}
For any $l\in\N$,
\begin{equation}\label{eq:decimated}
M^{2l} - \tr(M^l)\,M^l + \det(M^l)\,I \;=\; 0.
\end{equation}
Applying~\eqref{eq:decimated} to $\binom{x_k}{x_{k-1}}$ and taking the first component yields
\begin{equation}\label{eq:trisite}
x_{k+2l}-V_l\,x_{k+l}+Q^l\,x_k=0,\qquad V_l:=\tr(M^l),\ Q^l:=\det(M^l)=(\det M)^l.
\end{equation}
\end{proposition}
\begin{remark}[Lucas/Dickson polynomials]
If $\lambda_\pm$ are eigenvalues of $M$, then $V_l=\lambda_+^l+\lambda_-^l$ obeys $V_{n+1}=P V_n - Q V_{n-1}$ with $V_0=2$, $V_1=P$.
\end{remark}

\section{Events, anchor, and the Chebyshev--Lucas checksum}
\subsection{Eventization (operational)}\label{subsec:eventization}
Fix measurement cadence $\Delta t$. For each site $t$ let $E_t[n]$ be integrated power/flux in window $n$.
Choose a quantum/threshold $\varepsilon>0$ and set
\begin{equation}\label{eq:event}
\delta N_t[n]\;:=\;\Big\lfloor \frac{E_t[n]-E_t[n-1]}{\varepsilon}\Big\rfloor \in \Z_{\ge 0},\qquad
T_s:=\sum_n \delta N_s[n].
\end{equation}
An \emph{event} at $(t,n)$ is $\delta N_t[n]>0$. The \emph{anchor} experiment restricts events to $s=0$ (or pre--aligns them by virtual index shifts).

\subsection{Invariant and checksum; when ``symmetric regrouping'' applies}
Let $r$ be a root of $z^2 - V_l z + Q^l=0$ with $|r|\ge 1$ (the other root is $Q^l/r$).
For integer sequences $d=\{d_k\}$ with finite support define
\begin{equation}
I(d)\;=\;\sum_k d_k r^k.
\end{equation}
\begin{lemma}[Invariance under a tri--site move]\label{lem:invariance}
The local update $(d_{t-1},d_t,d_{t+1})\mapsto (d_{t-1}+Q^l,\ d_t-V_l,\ d_{t+1}+1)$ leaves $I(d)$ unchanged.
\end{lemma}
\begin{proof}
The change is $-V_l r^t + r^{t+1}+Q^l r^{t-1}=0$ by the characteristic equation.
\end{proof}

\begin{theorem}[Generalized checksum]\label{thm:checksum}
If $d$ is obtained from an event stream $\{\delta N_s[n]\}$ by any sequence of local tri--site moves, then
\begin{equation}\label{eq:general-checksum}
\sum_{k} d_k\,r^k \;=\; \sum_{s} T_s\, r^{s}.
\end{equation}
\end{theorem}
\begin{proof}
Each unit event at site $s$ increases $I$ by $r^s$. Lemma~\ref{lem:invariance} shows $I$ is invariant under subsequent tri--site moves.
\end{proof}

\paragraph{Symmetric regrouping (Dickson form).}
On a \emph{symmetric window} $[-R,R]$ and when the decoded digits satisfy $d_{-t}=d_t$ (e.g., due to symmetric injection/BC or explicit symmetrization of the window), one may regroup
\begin{equation}\label{eq:dickson}
\sum_{k=-R}^{R} d_k r^k \;=\; d_0 + \sum_{j=1}^{R} d_j\,\big(r^j + (Q^l/r)^j\big) \;=\; d_0 + \sum_{j=1}^{R} d_j\,D_j(V_l,Q^l),
\end{equation}
where $D_0=2$, $D_1=V_l$, and $D_{j+1}=V_l D_j - Q^l D_{j-1}$ are Dickson polynomials ($D_j(V_l,Q^l)=r^j+(Q^l/r)^j$).
Without symmetry, keep the two geometric progressions separately; the invariant~\eqref{eq:general-checksum} still holds.

\begin{example}[Illustration for $l=1$]\label{ex:l1}
Let $M=\begin{psmallmatrix}P&-Q\\1&0\end{psmallmatrix}$ with $P=1$, $Q=-1$ so $V_1=1$, $Q^1=-1$, and $r$ solves $r^2-r-1=0$ ($r=\varphi$).
Anchor events $T=5$ yield decoded digits $d_0=5$, others zero: $\sum_k d_k r^k=5$ matches~\eqref{eq:general-checksum}.
With $T=11$, a carry creates $(d_{-1},d_0,d_{+1})=(1,8,1)$, giving $11=8+1\cdot(r+(Q/r))=8+1\cdot D_1(V_1,Q^1)$.
\end{example}

\section{Decoding on a window: trellis, design rules, and streaming}
We observe $m_t=d_t+\eta_t$ on $t\in[-R,R]$.
\begin{definition}[Decoder (nearest--integer, constrained)]\label{def:decoder}
\begin{equation}\label{eq:decoder}
\hat d \in \arg\min_{d\in\mathcal D^{2R+1}} \ \sum_{t=-R}^{R} w_t\,|d_t - m_t|^2
\quad \text{s.t.}\quad d_{t+1}=V_l d_t - Q^l d_{t-1}\ \ \forall t\in[-R,R-1].
\end{equation}
\end{definition}
\begin{proposition}[Finite--state trellis]\label{prop:trellis}
States are pairs $(a,b)=(d_{t-1},d_t)\in\mathcal D^2$.
A transition $(a,b)\to (b,c)$ is allowed if $c=V_l b - Q^l a$ lies in $\mathcal D$ (soft trellis: $c=\mathrm{round}(\cdot)$ with a penalty).
Dynamic programming yields the ML path in $O(R\,|\mathcal D|^2)$ time.
\end{proposition}

\paragraph{Design rules (alphabet and window).}
Let $r>1$ be the larger root of $z^2-V_l z+Q^l=0$.
For nonnegative $\mathcal D^+=\{0,1,\dots,A\}$, avoid saturation by
\begin{equation}\label{eq:Amin}
A_{\min}(R)\ :=\ \Big\lceil \frac{T_{\max}}{\,1+\sum_{j=1}^{R} D_j(V_l,Q^l)\,} \Big\rceil,\qquad
A \;=\; \max\!\Big\{\,A_{\min}(R),\ \big\lfloor \sqrt{|V_l|}\,\big\rfloor \Big\},
\end{equation}
and choose
\begin{equation}\label{eq:Rchoice}
R \;=\; \Big\lceil \log_r\!\Big(1+\frac{(r-1)T_{\max}}{A}\Big)\Big\rceil - 1 \;+\; m,\qquad m\in\{2,3\}.
\end{equation}
Balanced $\mathcal D^\pm=\{-A,\dots,A\}$ uses the same rules.

\paragraph{Streaming (guard states).}
Maintain a sliding window. At each step, freeze the left guard pair $(d_{-R-1},d_{-R})$ from the previous solution,
prune transitions violating the recurrence at $t=-R$, and run one DP pass over $[-R{+}1,\dots,R{+}1]$.
This yields amortized $O(|\mathcal D|^2)$ updates per cadence.

\subsection{Noise tolerance (explicit constants)}\label{subsec:noise}
Assume $m_t=d_t+\eta_t$ with i.i.d.\ $\eta_t\sim\mathcal N(0,\sigma^2)$ and $d$ supported on $[-R,R]$.
A union bound over trellis paths gives
\begin{equation}\label{eq:sigmastar}
\mathbb P(\hat d \neq d) \ \le\ C_R \exp\!\Big(-\frac{d_{\min}^2}{8\sigma^2}\Big),
\qquad C_R \ \le\ (2R{+}1)\,|\mathcal D|^2,
\end{equation}
where $d_{\min}$ is the minimum Euclidean distance between distinct trellis paths and obeys the conservative estimate
\begin{equation}
d_{\min}^2 \ \gtrsim\ \min\{ r^{2R},\ (Q^l/r)^{2R}\}.
\end{equation}
Thus one admissible operating point is $\sigma_\star \le \frac{1}{\sqrt 8}\min\{ r^{R}, (Q^l/r)^{R}\}$.

\section{Support radius under repeated anchor injection}
Let $R(T)=\max\{|t|:\ d_t(T)\neq 0\}$ after $T$ anchor events and decoding.
Define
\begin{equation}\label{eq:SR}
S_R(V_l,Q^l)\ :=\ 1+\sum_{j=1}^{R} D_j(V_l,Q^l).
\end{equation}
\begin{proposition}[Maximal representable count at radius $R$]\label{prop:Tmax}
For $\mathcal D^+=\{0,\dots,A\}$,
\begin{equation}\label{eq:Tmax}
T_{\max}(R)\;=\; A\,S_R(V_l,Q^l).
\end{equation}
\end{proposition}
\begin{remark}[Why $S_R$ starts with $1$ (not $D_0$)]
The checksum~\eqref{eq:dickson} uses $d_0$ with coefficient $1$ (the term $r^0$), whereas $D_0=2$ corresponds to the symmetric combination $r^0 + (Q^l/r)^0=2$.
Maximizing $T$ at fixed radius sets all digits $(d_0,d_1,\dots,d_R)$ to $A$, yielding
$T_{\max}=A\,(1+\sum_{j=1}^{R} D_j)$, consistent with~\eqref{eq:SR}.
\end{remark}
Using $D_j(V_l,Q^l)=r^j+(Q^l/r)^j$ gives
\begin{equation}
\frac{r^{R+1}-1}{r-1}\ \le\ S_R(V_l,Q^l)\ \le\ \frac{2(r^{R+1}-1)}{r-1}.
\end{equation}
\begin{theorem}[Bounds and scaling]\label{thm:Rbounds}
\begin{equation}\label{eq:Rbounds}
\Big\lceil \log_r\!\Big(1+\frac{(r-1)T}{2A}\Big)\Big\rceil -1 \ \le\ R(T)\ \le\
\Big\lceil \log_r\!\Big(1+\frac{(r-1)T}{A}\Big)\Big\rceil -1.
\end{equation}
Hence $R(T)=\Theta(\log_r T)$ for $r>1$.
\end{theorem}
\begin{remark}[Marginal $r\approx 1$]
If $|V_l|\approx 2\sqrt{Q^l}$ then $r\approx 1$ and confinement degrades; choose the stride $l$ (supercell/Floquet) to enforce $|V_l|\ge 2\sqrt{Q^l}+\varepsilon$.
\end{remark}

\section{Uniqueness: what is and is not determined}\label{sec:uniqueness}
\paragraph{Checksum is not injective (counterexample).}
Let $\Delta$ be zero everywhere except at $(t{-}1,t,t{+}1)$ with values $(Q^l,\ -V_l,\ 1)$.
Then $\sum_k \Delta_k r^k=0$ by the characteristic equation. Thus the checksum \emph{alone} does not determine $d$.

\paragraph{Anchor case: uniqueness via normalization/ML on a window.}
On a finite window with a fixed alphabet and the recurrence constraint~\eqref{eq:decoder}, the ML decoder is a strictly convex quadratic over a finite trellis;
ties are broken deterministically by the DP backtrack, hence the decoded $\hat d$ is \emph{unique}.
This matches the canonical normal form selected by the local normalization calculus.

\paragraph{Multi--site events.}
For general $\{T_s\}$, the checksum~\eqref{eq:general-checksum} fixes $\sum_k d_k r^k$ but, by the counterexample, does not \emph{by itself} fix $d$.
Uniqueness is recovered by the same ML principle on a window (fixed alphabet, recurrence constraints) and, if desired, symmetric windows (which eliminate sign--flip ambiguities by construction).
When per--site totals are known, one may pre--align to the anchor frame and decode per shift; the superposition is then validated by the generalized checksum.

\section{Cross--track coupling: per--mode generalization and norms}
For a multi--track chain with nearest--track coupling, transverse Fourier diagonalization yields per--mode transfers $M(\theta)$.
We use the operator norm $\|\cdot\|_2$ unless otherwise stated.
\begin{proposition}[Per--mode decimation]\label{prop:modal}
Each mode $\theta$ obeys
\begin{equation}\label{eq:mode}
x^{(\theta)}_{k+2l} - V_l(\theta)\,x^{(\theta)}_{k+l} + Q^l(\theta)\,x^{(\theta)}_k=0,\qquad
V_l(\theta)=\tr(M(\theta)^l),\ \ Q^l(\theta)=\det(M(\theta)^l).
\end{equation}
\end{proposition}
\paragraph{Quantitative robustness.}
If $M(\theta)\mapsto M(\theta)+E(\theta)$ with perturbation $E$, then
\begin{equation}
\big|\tr((M+E)^l)-\tr(M^l)\big| \ \le\ l\,\|M\|_2^{\,l-1}\,\|E\|_2 + O(\|E\|_2^2).
\end{equation}
We call coupling \emph{weak at level $\delta$} if $l\,\|M\|_2^{\,l-1}\,\|E\|_2\le \delta\,|\tr(M^l)|$ uniformly in $\theta$;
under this condition the per--mode checksum is accurate to $O(\delta)$.

\section{Choosing the stride $l$ (natural origins)}
\begin{itemize}
\item \textbf{Supercell:} $l$ equals the microperiod (geometry).
\item \textbf{Floquet locking:} with phase advance $\theta$ per site, pick $l$ s.t.\ $\theta/\pi=p/l$ rational, making $V_l=2\cos(l\theta)$ typically small integers.
\item \textbf{Topology:} on a ring, one round--trip per residue class $\Rightarrow$ decimation by $l$; APBC ($\pi$--flux) flips the global holonomy.
\end{itemize}

\section{Experimental protocols and validation metrics}
P0--P4 as in the main version, with added plots: checksum accuracy vs.\ $\sigma/\sigma_\star$, decoder success rate vs.\ $\sigma/\sigma_\star$, and $R(T)$ scaling with the theoretical bounds.

\section{Related work (selected)}
\begin{itemize}
\item \textbf{Transfer matrices and decimation:} G.\ Economou, \emph{Green's Functions in Quantum Physics}, 3rd ed., Springer, 2006.
\item \textbf{Chebyshev polynomials:} T.\ J.\ Rivlin, \emph{Chebyshev Polynomials}, 2nd ed., Wiley, 1990.
\item \textbf{Dickson polynomials:} R.\ Lidl and H.\ Niederreiter, \emph{Finite Fields}, Cambridge Univ.\ Press, 1997 (see chs.\ on Dickson polynomials).
\item \textbf{Convolutional codes and trellises:} A.\ J.\ Viterbi, ``Error bounds for convolutional codes and an asymptotically optimum decoding algorithm,'' \emph{IEEE Trans.\ Inf.\ Theory}, 1967; G.\ D.\ Forney Jr., ``Convolutional codes I: Algebraic structure,'' \emph{IEEE Trans.\ Inf.\ Theory}, 1970.
\end{itemize}
\paragraph{Contrast to classical trellises.}
Classical trellises are engineered algebraically; here the parity constraint (kernel) and alphabet are \emph{medium--determined} via $(V_l,Q^l)$ from the unit cell and the chosen stride $l$ (supercell/Floquet/topology). This enables error--checked, event--coded metrology without explicit coding hardware.

\section*{Acknowledgments}
We thank the transfer--matrix, special polynomial, and coding theory communities for foundational ideas.

\appendix
\section{Additional proofs}
\begin{proof}[Proof of Proposition~\ref{prop:decimated}]
Cayley--Hamilton on $M^l$ gives $M^{2l}-(\tr M^l)M^l+(\det M^l)I=0$. Acting on $\binom{x_k}{x_{k-1}}$ and taking the first component yields~\eqref{eq:trisite}.
\end{proof}
\begin{proof}[Lemma~\ref{lem:invariance}]
Immediate from $r^2-V_l r + Q^l=0$.
\end{proof}
\end{document}
